% =============================================================================
% 5.2 USER FLOWS (FLUJOS DE NAVEGACIÓN DEL PERFIL DE USUARIO)
% =============================================================================
\section{User Flows (Flujos de Navegación del Perfil de Usuario)}
\label{sec:user_flows}

Los \textit{user flows} modelan la experiencia de navegación de los dos perfiles
principales del sistema, trazando el camino desde la entrada hasta el objetivo de valor de
cada rol. Estos flujos materializan los casos de uso definidos en el
capítulo~\ref{ch:requisitos} (pág.~\pageref{ch:requisitos}).

\subsection{Registro, configuración inicial y carga de portafolio}
\label{subsec:flow_registro}

La figura~\ref{fig:flow-onboarding} representa el embudo de incorporación de un
profesional: desde el registro hasta la publicación de su portafolio, incluyendo la
ramificación por rol que ocurre inmediatamente tras el alta.

\vspace{0.3cm}
\begin{figure}[h!]
\centering
\begin{tikzpicture}[node distance=0.8cm and 1.1cm]
    \node[flujoInicio] (start) {Registro};
    \node[flujoDecision, below=of start] (rol) {¿Rol?};
    \node[flujoProceso, below left=0.9cm and 0.2cm of rol] (prof) {Dashboard\\Profesional};
    \node[flujoProceso, below right=0.9cm and 0.2cm of rol] (recl) {Dashboard\\Reclutador};
    \node[flujoProceso, below=of prof] (perfil) {Configurar perfil\\(bio, foto, skills)};
    \node[flujoProceso, below=of perfil] (proy) {Cargar proyectos\\y experiencia};
    \node[flujoProceso, below=of proy] (ajustes) {Ajustes de\\portafolio (SEO)};
    \node[flujoInicio, fill=c4Sistema, below=of ajustes] (pub) {Portafolio público};

    \draw[flujoFlecha] (start) -- (rol);
    \draw[flujoFlecha] (rol) -- node[c4etiqueta, left]{PROFESSIONAL} (prof);
    \draw[flujoFlecha] (rol) -- node[c4etiqueta, right]{RECRUITER} (recl);
    \draw[flujoFlecha] (prof) -- (perfil);
    \draw[flujoFlecha] (perfil) -- (proy);
    \draw[flujoFlecha] (proy) -- (ajustes);
    \draw[flujoFlecha] (ajustes) -- (pub);
\end{tikzpicture}
\caption{Flujo de incorporación del profesional: del registro al portafolio público.}
\label{fig:flow-onboarding}
\end{figure}

El paso de \textbf{configurar experiencia} incorpora una sub-experiencia geolocalizada:
en \pathbreak{/dashboard/experience} el usuario sitúa cada empleo en un mapa interactivo
(Leaflet) y el sistema deriva la dirección textual por geocodificación inversa, como se
detalla en la sección~\ref{subsec:comp_avanzados} (pág.~\pageref{subsec:comp_avanzados}).

\subsection{Embudo del reclutador (búsqueda, filtrado y analíticas)}
\label{subsec:flow_reclutador}

El reclutador opera un \textbf{embudo de descubrimiento}: parte de una búsqueda amplia
(filtros estructurados o lenguaje natural traducido por IA), refina hacia un
\textit{pipeline} gestionable mediante un tablero Kanban, y culmina en analíticas e
\textit{insights} generados por IA. La figura~\ref{fig:flow-recruiter} lo ilustra.

\vspace{0.3cm}
\begin{figure}[h!]
\centering
\begin{tikzpicture}[node distance=0.75cm and 1.4cm]
    \node[flujoInicio] (start) {Talent Discovery};
    \node[flujoProceso, below=of start] (search) {Buscar (filtros / NLP)};
    \node[flujoProceso, below=of search] (kanban) {Organizar en Kanban\\(drag \& drop)};
    \node[flujoProceso, below=of kanban] (notas) {Notas privadas\\+ match ring};
    \node[flujoProceso, below=of notas] (dash) {Dashboard\\(funnel · radar · KPIs)};
    \node[flujoInicio, fill=c4Sistema, below=of dash] (ai) {Insights IA + chat};

    \draw[flujoFlecha] (start) -- (search);
    \draw[flujoFlecha] (search) -- (kanban);
    \draw[flujoFlecha] (kanban) -- (notas);
    \draw[flujoFlecha] (notas) -- (dash);
    \draw[flujoFlecha] (dash) -- (ai);
\end{tikzpicture}
\caption{Embudo de descubrimiento de talento del reclutador.}
\label{fig:flow-recruiter}
\end{figure}

\begin{AvisoNota}
El paso \textbf{Buscar (NLP)} invoca el modo \texttt{nlp\_search} del motor de IA
(capítulo~\ref{ch:backend}, tabla~\ref{tab:ia-modos}, pág.~\pageref{tab:ia-modos}), que
traduce una frase como ``desarrolladores React senior en Bolivia'' a un objeto de filtros
estructurados que el frontend aplica sobre el catálogo de talento. La IA no consulta la
base de datos: solo estructura la intención del reclutador.
\end{AvisoNota}

\subsection{Visitante anónimo de un portafolio público}
\label{subsec:flow_visitante}

El tercer flujo relevante es el del \textbf{visitante no autenticado} que consulta un
portafolio compartido (\pathbreak{/p/:slug}). Este recorrido no requiere sesión y puede
estar protegido por contraseña según los ajustes del dueño. La
figura~\ref{fig:flow-visitante} lo modela.

\vspace{0.3cm}
\begin{figure}[h!]
\centering
\begin{tikzpicture}[node distance=0.75cm and 1.1cm]
    \node[flujoInicio] (start) {Visita \pathbreak{/p/:slug}};
    \node[flujoDecision, below=of start] (pwd) {¿Protegido\\por clave?};
    \node[flujoProceso, right=2.4cm of pwd] (gate) {Pantalla de\\contraseña};
    \node[flujoDecision, below=1.0cm of pwd] (pub) {¿Publicado?};
    \node[flujoProceso, right=2.4cm of pub] (404) {No encontrado\\(404)};
    \node[flujoInicio, fill=c4Sistema, below=1.0cm of pub] (view) {Portafolio renderizado};

    \draw[flujoFlecha] (start) -- (pwd);
    \draw[flujoFlecha] (pwd) -- node[c4etiqueta, above]{sí} (gate);
    \draw[flujoFlecha] (gate) |- node[c4etiqueta, near start]{ok} (pub);
    \draw[flujoFlecha] (pwd) -- node[c4etiqueta, left]{no} (pub);
    \draw[flujoFlecha] (pub) -- node[c4etiqueta, above]{no} (404);
    \draw[flujoFlecha] (pub) -- node[c4etiqueta, left]{sí} (view);
\end{tikzpicture}
\caption{Flujo del visitante anónimo a un portafolio público.}
\label{fig:flow-visitante}
\end{figure}

Este flujo consume exclusivamente \textit{endpoints} públicos
(\pathbreak{/api/v1/portfolio/public/**}) declarados en \texttt{PUBLIC\_PATHS}
(capítulo~\ref{ch:backend}, tabla~\ref{tab:public-paths}, pág.~\pageref{tab:public-paths}),
sin token; las políticas RLS del motor garantizan que solo se exponga contenido marcado
como público y publicado.
